\documentclass[12pt]{article}
\usepackage[utf8]{inputenc}
\usepackage[T1]{fontenc}
\usepackage{hyperref}
\title{Yalakalore Language Resources for K-12 Curriculum Development}
\author{Compiled by Claude Sonnet 5.5 (claude-sonnet-5-5) for Andr\'as Kornai}
\date{}
\begin{document}
\maketitle

\section{Introduction}
Yalakalore (ISO 639-3: xyl; Glottocode yala1242) is an extinct language of Brazil that Glottolog lists as unattested. It is reported in connection with the Nambikwara family of Mato Grosso, with Alapmunte, Tawend\^e and Yelelihre listed alongside it. Searches found no recorded words or texts.

\section{Primary Educational Resources}
None exist as far as the searches showed.

\section{Linguistic and Reference Materials}
\begin{itemize}
\item Glottolog record for ISO code xyl (Yalakalore). \url{https://glottolog.org/resource/languoid/iso/xyl}
\item David M. Eberhard (2009), overview of the Nambikwara family listing its varieties (Glottolog reference record). \url{https://glottolog.org/resource/reference/id/37160}
\item Wikipedia, Nambikwaran languages (family context). \url{https://en.wikipedia.org/wiki/Nambikwaran_languages}
\end{itemize}

\section{Religious and Community Resources}
None found.

\section{Digital and Audio Resources}
None found.

\section{Curriculum Development Context}
Because the language is unattested, no curriculum can teach it. A school could instead teach about the living Nambikwara languages and peoples of Mato Grosso, using materials produced by those communities, and could explain how languages come to be known only by name.
\end{document}